\documentclass[11pt]{article}
\usepackage{amsmath,amssymb,amsthm}
\usepackage{hyperref}
\usepackage{algorithm}
\usepackage{algorithmic}
\usepackage{tikz}
\usetikzlibrary{arrows.meta,positioning}

% Theorem environments
\newtheorem{definition}{Definition}
\newtheorem{theorem}{Theorem}
\newtheorem{lemma}{Lemma}
\newtheorem{proposition}{Proposition}
\newtheorem{corollary}{Corollary}

\title{OBINexus Cryptographic Substrate:\\Formal Specification with Dimensional Game Theory}
\author{OBINexus Consortium}
\date{\today}

\begin{document}
\maketitle

\begin{abstract}
This document provides a formal mathematical specification of the OBINexus cryptographic primitive substrate, incorporating dimensional game theory for key derivation and rotation strategies. We establish the security properties of the HKDF-based rehash function within a multi-dimensional strategic framework where cryptographic dimensions represent distinct security domains.
\end{abstract}

\section{Introduction}

The OBINexus cryptographic substrate operates on the principle of dimensional game theory, where each cryptographic operation exists within a strategic dimension $\mathcal{D}_i$. The ``Sue Sue Sue'' trilogy manifests as three fundamental dimensions: Security ($\mathcal{D}_S$), Stability ($\mathcal{D}_T$), and Sovereignty ($\mathcal{D}_V$).

\section{Dimensional Game Theory Framework}

\subsection{Formal Definition of Dimensions}

\begin{definition}[Cryptographic Dimension]
A cryptographic dimension $\mathcal{D}$ is a tuple $(\mathcal{S}, \mathcal{A}, \mathcal{P}, u)$ where:
\begin{itemize}
    \item $\mathcal{S}$ is the state space of cryptographic primitives
    \item $\mathcal{A}$ is the action space of operations
    \item $\mathcal{P}: \mathcal{S} \times \mathcal{A} \rightarrow \mathcal{S}$ is the transition function
    \item $u: \mathcal{S} \rightarrow \mathbb{R}$ is the security utility function
\end{itemize}
\end{definition}

\begin{definition}[Dimensional Game]
A dimensional game $\Gamma$ is defined as:
\[\Gamma = \langle N, \{\mathcal{D}_i\}_{i \in N}, \{\sigma_i\}_{i \in N}, \{\pi_i\}_{i \in N} \rangle\]
where:
\begin{itemize}
    \item $N = \{1, 2, ..., n\}$ is the set of players (system components)
    \item $\mathcal{D}_i$ is the dimension for player $i$
    \item $\sigma_i: \mathcal{S}_i \rightarrow \Delta(\mathcal{A}_i)$ is player $i$'s strategy
    \item $\pi_i: \prod_{j \in N} \mathcal{S}_j \rightarrow \mathbb{R}$ is player $i$'s payoff function
\end{itemize}
\end{definition}

\subsection{The Three Dimensions}

\begin{definition}[Security Dimension $\mathcal{D}_S$]
The security dimension is defined as:
\[\mathcal{D}_S = (\mathcal{S}_S, \mathcal{A}_S, \mathcal{P}_S, u_S)\]
where:
\begin{align}
\mathcal{S}_S &= \{0, 1\}^{256} \times \{0, 1\}^{128} \times \mathbb{N} \quad \text{(keys × salts × counters)}\\
\mathcal{A}_S &= \{\text{derive}, \text{rotate}, \text{extract}, \text{expand}\}\\
u_S(s) &= -H(K|C) \quad \text{(negative conditional entropy)}
\end{align}
\end{definition}

\begin{definition}[Stability Dimension $\mathcal{D}_T$]
\[\mathcal{D}_T = (\mathcal{S}_T, \mathcal{A}_T, \mathcal{P}_T, u_T)\]
where $u_T$ measures collision resistance and fault tolerance.
\end{definition}

\begin{definition}[Sovereignty Dimension $\mathcal{D}_V$]
\[\mathcal{D}_V = (\mathcal{S}_V, \mathcal{A}_V, \mathcal{P}_V, u_V)\]
where $u_V$ quantifies control over key exposure and rotation autonomy.
\end{definition}

\section{HKDF-Based Rehash Function}

\subsection{Formal Specification}

\begin{definition}[Rehash Function]
The rehash function $\mathcal{R}: \{0,1\}^* \times \{0,1\}^* \times \{0,1\}^* \rightarrow \{0,1\}^{256}$ is defined as:
\[\mathcal{R}(\text{salt}, K_{pub}, K_{priv}) = \text{HKDF}(\text{salt}, K_{priv}, K_{pub}, 256)\]
where HKDF is composed of:
\begin{align}
\text{PRK} &= \text{HKDF-Extract}(\text{salt}, K_{priv})\\
\text{OKM} &= \text{HKDF-Expand}(\text{PRK}, K_{pub}, 256)
\end{align}
\end{definition}

\begin{theorem}[Security of Rehash]
Given that HMAC-SHA256 is a pseudorandom function (PRF), the rehash function $\mathcal{R}$ satisfies:
\[\Pr[\mathcal{A}^{\mathcal{R}(\cdot,\cdot,K_{priv})}(1^{\lambda}) = 1] - \Pr[\mathcal{A}^{f(\cdot,\cdot)}(1^{\lambda}) = 1] \leq \text{negl}(\lambda)\]
for any PPT adversary $\mathcal{A}$ and random function $f: \{0,1\}^* \times \{0,1\}^* \rightarrow \{0,1\}^{256}$.
\end{theorem}

\subsection{Dimensional Properties}

\begin{proposition}[Cross-Dimensional Security]
For any strategy profile $(\sigma_S, \sigma_T, \sigma_V)$ in the dimensional game $\Gamma$:
\[\mathbb{E}[u_S + u_T + u_V] \geq \sum_{i \in \{S,T,V\}} \mathbb{E}[u_i | \sigma_i \text{ alone}]\]
This shows that coordinated strategies across dimensions yield superior security.
\end{proposition}

\section{Substrate State Transitions}

\subsection{State Evolution}

The substrate state evolves according to:
\[\mathcal{S}_{t+1} = \mathcal{P}(\mathcal{S}_t, a_t)\]
where the transition function $\mathcal{P}$ is defined as:

\begin{algorithm}
\caption{Substrate State Transition}
\begin{algorithmic}
\STATE \textbf{Input:} Current state $s_t$, action $a$
\STATE \textbf{Output:} New state $s_{t+1}$
\IF{$a = \text{rotate}$}
    \STATE $\text{salt}_{new} \leftarrow \{0,1\}^{128}$ (uniform random)
    \STATE $K_{new} \leftarrow \mathcal{R}(\text{salt}_{new}, K_{pub}, K_{priv})$
    \STATE $c_{t+1} \leftarrow c_t + 1$ (increment counter)
\ELSIF{$a = \text{expose}$}
    \STATE Apply exposure control $\mathcal{E}(K, \text{level})$
\ENDIF
\RETURN $s_{t+1}$
\end{algorithmic}
\end{algorithm}

\subsection{Equilibrium Analysis}

\begin{theorem}[Nash Equilibrium in Dimensional Game]
There exists a Nash equilibrium $(\sigma_S^*, \sigma_T^*, \sigma_V^*)$ where:
\begin{itemize}
    \item $\sigma_S^*$: Rotate keys with period $\tau^* = \arg\max_{\tau} \frac{u_S(\tau)}{c(\tau)}$
    \item $\sigma_T^*$: Use atomic updates for all critical operations
    \item $\sigma_V^*$: Apply controlled exposure with zero-knowledge proofs
\end{itemize}
\end{theorem}

\section{Security Analysis}

\subsection{Collision Resistance}

\begin{lemma}[Birthday Bound]
For the rehash function $\mathcal{R}$ with 256-bit output, the probability of collision after $q$ queries is:
\[\Pr[\text{collision}] \leq \frac{q^2}{2^{257}}\]
\end{lemma}

\subsection{Side-Channel Resistance}

\begin{definition}[Timing-Safe Implementation]
An implementation is $\epsilon$-timing-safe if:
\[\max_{x,y} |T(x) - T(y)| \leq \epsilon\]
where $T(x)$ is the execution time on input $x$.
\end{definition}

\begin{proposition}[HMAC Timing Properties]
The HMAC construction in HKDF provides constant-time operation when:
\begin{enumerate}
    \item Key comparison uses constant-time algorithms
    \item Hash function processes blocks in constant time
    \item No data-dependent branching occurs
\end{enumerate}
\end{proposition}

\section{Formal Verification Properties}

\subsection{Invariants}

The following invariants must hold for all substrate operations:

\begin{align}
\text{INV}_1 &: \forall t, \mathcal{H}(K_t) \geq \lambda \quad \text{(minimum entropy)}\\
\text{INV}_2 &: \forall i \neq j, \Pr[K_i = K_j] \leq 2^{-256} \quad \text{(uniqueness)}\\
\text{INV}_3 &: \forall t, c_{t+1} > c_t \quad \text{(monotonic counter)}
\end{align}

\subsection{Dimensional Coupling}

\begin{theorem}[Dimensional Coupling Theorem]
The security of the overall system is lower-bounded by:
\[\text{Sec}(\Gamma) \geq \min_{i \in \{S,T,V\}} \text{Sec}(\mathcal{D}_i) - \epsilon_{coupling}\]
where $\epsilon_{coupling} = O(2^{-\lambda})$ represents the coupling loss between dimensions.
\end{theorem}

\section{Implementation Constraints}

\subsection{Substrate Layers}

\begin{center}
\begin{tikzpicture}[
    node distance=2cm,
    every node/.style={draw,rectangle,align=center,text width=3cm}
]
\node (s0) {Substrate 0\\Max Security};
\node[below of=s0] (s1) {Substrate 1\\Staged Deploy};
\node[below of=s1] (s2) {Substrate 2\\Application};

\draw[->] (s2) -- (s1) node[midway,right,draw=none,text width=auto] {promote};
\draw[->] (s1) -- (s0) node[midway,right,draw=none,text width=auto] {promote};
\draw[<-,dashed] (s2) -- (s1) node[midway,left,draw=none,text width=auto] {expose};
\draw[<-,dashed] (s1) -- (s0) node[midway,left,draw=none,text width=auto] {expose};
\end{tikzpicture}
\end{center}

\subsection{Performance Bounds}

\begin{proposition}[Computational Complexity]
The rehash function has complexity:
\begin{itemize}
    \item Time: $O(\ell)$ where $\ell = |K_{priv}| + |K_{pub}| + |\text{salt}|$
    \item Space: $O(1)$ (constant memory)
    \item Rounds: $\lceil \frac{256}{256} \rceil = 1$ (for 256-bit output)
\end{itemize}
\end{proposition}

\section{Quantum-Resistant Extensions}

\subsection{Quantum Dimensional Game Theory}

\begin{definition}[Quantum Cryptographic Dimension]
A quantum dimension $\mathcal{Q}_i$ extends the classical dimension with quantum states:
\[\mathcal{Q}_i = (\mathcal{H}_i, \mathcal{U}_i, \mathcal{M}_i, \rho_i)\]
where:
\begin{itemize}
    \item $\mathcal{H}_i$ is the Hilbert space of quantum states
    \item $\mathcal{U}_i$ is the set of unitary operations
    \item $\mathcal{M}_i$ is the set of quantum measurements
    \item $\rho_i \in \mathcal{D}(\mathcal{H}_i)$ is the density matrix
\end{itemize}
\end{definition}

\subsection{QFT-Based Security Foundation}

\begin{definition}[Quantum Fourier Transform Security Layer]
The QFT security layer implements:
\[\text{QFT}_{N}|j\rangle = \frac{1}{\sqrt{N}} \sum_{k=0}^{N-1} e^{2\pi ijk/N}|k\rangle\]
Applied to cryptographic primitives:
\[\mathcal{K}_{quantum} = \text{QFT} \circ \mathcal{R} \circ \text{QFT}^{\dagger}\]
\end{definition}

\begin{theorem}[Quantum Security Against Grover]
Against Grover's algorithm, the rehash function with quantum-resistant modifications achieves:
\[\text{Security}_{quantum} \geq 2^{128} \text{ quantum operations}\]
by using 256-bit classical security parameters.
\end{theorem}

\subsection{Post-Quantum Rehash Function}

\begin{definition}[PQ-Rehash]
The post-quantum rehash extends classical rehash with lattice-based hardness:
\begin{align}
\mathcal{R}_{PQ}(\text{salt}, K_{pub}, K_{priv}) &= \mathcal{R}(\text{salt}, K_{pub}, K_{priv}) \oplus \mathcal{L}(K_{priv})\\
\mathcal{L}(K) &= \text{LWE}_{n,q,\chi}(K) \mod 2^{256}
\end{align}
where LWE is the Learning With Errors problem with dimension $n$, modulus $q$, and error distribution $\chi$.
\end{definition}

\subsection{Quantum-Resistant Substrate Layers}

\begin{algorithm}
\caption{Quantum-Secure Substrate Update}
\begin{algorithmic}
\STATE \textbf{Input:} Classical state $s_t$, quantum state $|\psi\rangle$
\STATE \textbf{Output:} Updated hybrid state $(s_{t+1}, |\psi'\rangle)$
\STATE // Anti-commercialization: Real security, not marketing
\IF{$\text{QuantumThreat}(s_t) > \tau_{quantum}$}
    \STATE $K_{lattice} \leftarrow \text{SampleLatticeKey}(2048)$
    \STATE $K_{hash} \leftarrow \text{SPHINCS+}(K_{priv})$
    \STATE $K_{hybrid} \leftarrow K_{classical} \oplus K_{lattice} \oplus K_{hash}$
    \STATE $|\psi'\rangle \leftarrow \text{QFT}(|K_{hybrid}\rangle)$
\ELSE
    \STATE Use classical update with quantum monitoring
\ENDIF
\RETURN $(s_{t+1}, |\psi'\rangle)$
\end{algorithmic}
\end{algorithm}

\subsection{Security Against Quantum Adversaries}

\begin{theorem}[Post-Quantum Dimensional Security]
For a quantum adversary $\mathcal{A}_{quantum}$ with polynomial quantum gates:
\[\Pr[\mathcal{A}_{quantum} \text{ breaks } \mathcal{R}_{PQ}] \leq \text{negl}(\lambda) + \epsilon_{LWE} + \epsilon_{SPHINCS}\]
where:
\begin{itemize}
    \item $\epsilon_{LWE} = 2^{-\Omega(n/\log n)}$ for LWE parameters
    \item $\epsilon_{SPHINCS} = 2^{-256}$ for SPHINCS+ security
\end{itemize}
\end{theorem}

\begin{proposition}[Money vs Security Principle]
While quantum computing commercialization optimizes for $\max(\text{profit})$, OBINexus optimizes for:
\[\max(\text{Security}_{real}) \text{ subject to } \text{Cost} \leq \text{Cost}_{classical} \times 1.5\]
This ensures genuine security without quantum hype inflation.
\end{proposition}

\subsection{Quantum Dimensional Equilibrium}

\begin{theorem}[Quantum Nash Equilibrium]
In the quantum dimensional game, the equilibrium strategy includes:
\begin{enumerate}
    \item Classical dimensions: Continue HKDF-based operations
    \item Quantum dimension: Layer post-quantum primitives when $P(\text{quantum threat}) > 0.01$
    \item Hybrid mode: $\sigma_{hybrid}^* = \alpha\sigma_{classical} + (1-\alpha)\sigma_{quantum}$
\end{enumerate}
where $\alpha = e^{-t/\tau_{quantum}}$ decreases over time.
\end{theorem}

\subsection{Implementation Without Hype}

\begin{definition}[No-Nonsense Quantum Security]
Real quantum resistance requires:
\begin{itemize}
    \item Lattice keys: $n \geq 1024$, $q \approx 2^{23}$
    \item Hash-based signatures: SPHINCS+ with 256-bit security
    \item No reliance on unproven "quantum-safe" marketing claims
    \item Hybrid approach maintaining classical security
\end{itemize}
\end{definition}

\section{Conclusion}

The OBINexus cryptographic substrate provides a formally verified foundation for secure key derivation and rotation within a dimensional game theory framework. The ``Sue Sue Sue'' philosophy is mathematically instantiated through the three-dimensional strategic space, ensuring security, stability, and sovereignty properties are maintained throughout the system lifecycle.

With quantum-resistant extensions, the system achieves genuine post-quantum security without falling prey to commercialization hype. The QFT foundation and hybrid classical-quantum approach ensure practical security against both current and future threats.

\subsection{Future Work}

\begin{enumerate}
    \item Implement lattice-based dimensional transitions
    \item Formalize quantum multi-party games with entanglement
    \item Prove zero-knowledge properties for quantum substrate states
    \item Develop cost-effective quantum security without marketing inflation
\end{enumerate}

\end{document}